% Источник: лекция 7, с. 1--4; лекция 8, с. 8;
% лекция 9, с. 11--14; лекция 10, с. 18 рукописного PDF.

\section{Случайные величины}

\begin{flushright}
    \textit{Эмоции = результат $-$ ожидания.}
\end{flushright}

Пусть $(\Omega,\mathcal{F},\Pr)$ --- вероятностное пространство.
Идея: случайная величина --- числовая характеристика
случайного эксперимента.

\begin{definition}
    Функция $\xi\colon\Omega\to\R$ называется
    \textbf{случайной величиной}, если
    \[
        \xi^{-1}\bigl((-\infty;x]\bigr)\in\mathcal{F}
        \qquad\forall x\in\R.
    \]
    Эквивалентная запись:
    \[
        \{\omega\in\Omega\colon\xi(\omega)\leq x\}
        \in\mathcal{F}
        \qquad\forall x\in\R.
    \]
    Эквивалентное условие:
    \[
        \xi^{-1}(B)\in\mathcal{F}
        \qquad\forall B\in\mathcal{B}(\R).
    \]
    Здесь $\mathcal{B}(\R)$ --- наименьшая $\sigma$-алгебра,
    содержащая промежутки на $\R$, в частности
    $(y;x]$, $(y;x)$, $[y;x]$, $[y;x)$.
\end{definition}

\begingroup
\tcbset{myboxstyle/.append style={breakable=false}}
\begin{example}[Температура]
    Пусть $t(\omega)$ --- температура в градусах Цельсия.
    Рассматриваем вероятность \mbox{$\Pr\{t(\omega)>25\}$}.
    Соответствующее событие:
    \[
        \begin{aligned}
            \{\omega\in\Omega\colon t(\omega)>25\}
            &=t^{-1}\bigl((25;+\infty)\bigr)\\
            &=\Omega\setminus t^{-1}\bigl((-\infty;25]\bigr)
            \in\mathcal{F}.
        \end{aligned}
    \]
\end{example}
\endgroup

\subsection{Распределение случайной величины}

Идея: распределение --- то, какие значения и с какими
вероятностями принимает случайная величина.

\begin{definition}
    \textbf{Распределением случайной величины} $\xi$ называется
    вероятностная мера на $\mathcal{B}(\R)$, определённая равенством
    \[
        \Pr_\xi(B)=\Pr\{\xi\in B\},
        \qquad B\in\mathcal{B}(\R).
    \]
\end{definition}

\begin{center}
    \begin{tikzpicture}[>=Stealth]
        \draw[thick] (0,0) ellipse (0.95cm and 1.15cm);
        \node at (0,0) {$\Omega$};
        \node at (0,-1.55) {$(\Omega,\mathcal{F},\Pr)$};
        \draw[->,thick] (1.25,0) -- node[above] {$\xi$} (4.05,0);
        \draw[->,thick] (4.45,-1.2) -- (4.45,1.35);
        \node[right] at (4.45,1.15) {$\R$};
        \node[right] at (4.45,0.6) {$\mathcal{B}(\R)$};
        \node[right] at (4.45,-0.8)
            {$\Pr_\xi(B)=\Pr\{\xi\in B\}$};
        \draw[->] (0.55,-1.8) to[bend right=15] (4.35,-1.15);
    \end{tikzpicture}
\end{center}

\subsection{Дискретные распределения}

\begin{definition}
    \textbf{Дискретным распределением} называется распределение
    случайной величины, принимающей не более чем счётное
    число значений.
\end{definition}

Для дискретного распределения $\{(x_i;p_i)\}_i$:
\[
    \Pr_\xi(B)=\sum_{j\colon x_j\in B}p_j,
    \qquad B\in\mathcal{B}(\R).
\]

\begin{example}[Сумма очков двух кубиков]
    \[
        \omega=(i_1,i_2),\qquad i_j\in\{1,\ldots,6\},
        \qquad \xi(\omega)=i_1+i_2.
    \]
    \begin{center}
        \renewcommand{\arraystretch}{1.7}
        \setlength{\tabcolsep}{5pt}
        \begin{tabular}{c|*{11}{c}}
            $x_i$ & $2$ & $3$ & $4$ & $5$ & $6$ & $7$
                  & $8$ & $9$ & $10$ & $11$ & $12$ \\
            \hline
            $p_i$ & $\frac1{36}$ & $\frac2{36}$ & $\frac3{36}$
                  & $\frac4{36}$ & $\frac5{36}$ & $\frac6{36}$
                  & $\frac5{36}$ & $\frac4{36}$ & $\frac3{36}$
                  & $\frac2{36}$ & $\frac1{36}$
        \end{tabular}
    \end{center}
    \[
        \Pr\{\xi(\omega)=2\}=\frac1{36}.
    \]
\end{example}

\subsubsection{Бернуллиевское распределение}

\begin{definition}
    Смысл: Результат одного испытания: 1 — успех, 0 — неудача
    \[
        \begin{array}{c|cc}
            x_i & 0 & 1\\\hline
            p_i & 1-p & p
        \end{array}
    \]
\end{definition}

\subsubsection{Биномиальное распределение}

\begin{definition}
    \[
        \operatorname{Bin}(n;p),\qquad n\in\N,\quad p\in(0;1).
    \]
    \[
        p_k=\Pr\{\xi(\omega)=k\}=C_n^k p^k(1-p)^{n-k},
        \qquad k\in\{0,\ldots,n\}.
    \]
\end{definition}

Смысл: число успехов в схеме испытаний Бернулли.
\[
    \omega=(i_1,\ldots,i_n),\qquad i_j\in\{0;1\},
    \qquad \xi(\omega)=i_1+\cdots+i_n.
\]
При $k\in\{0,\ldots,n\}$:
\[
    \begin{aligned}
        p_k&=\Pr\{\xi(\omega)=k\}
        =\sum_{\omega\colon\sum_{j=1}^{n}i_j=k}\Pr(\omega)\\
        &=\sum_{\omega\colon\sum_{j=1}^{n}i_j=k}
          p^k(1-p)^{n-k}
        =C_n^k p^k(1-p)^{n-k}.
    \end{aligned}
\]
\[
    \sum_{k=0}^{n}C_n^k p^k(1-p)^{n-k}
    =\bigl(p+(1-p)\bigr)^n=1.
\]

\begin{example}
    При $n=10^4$, $p=0{,}003$:
    \[
        \Pr\{\xi\leq5\}
        =\sum_{k=0}^{5}C_{10^4}^k\,
        0{,}003^k\cdot0{,}997^{10^4-k}.
    \]
\end{example}

\subsubsection{Пуассоновское распределение}

\begin{definition}
    \[
        \operatorname{Poiss}(\lambda),\qquad\lambda>0.
    \]
    \[
        x_k=k\in\{0,1,2,\ldots\},\qquad
        p_k=\frac{\lambda^k}{k!}e^{-\lambda}.
    \]
\end{definition}

\begingroup
\small
\setlength{\abovedisplayskip}{6pt}
\setlength{\belowdisplayskip}{6pt}
\noindent Пояснение.
Здесь $p_k=\Pr\{\xi=k\}$ --- вероятность значения $k$.
Сумма $p_0+p_1+\cdots$ всех вероятностей должна равняться $1$.
Используем ряд экспоненты:
$\textstyle\sum_{k=0}^{\infty}\lambda^k/k!=e^\lambda$.
\[
    \sum_{k=0}^{\infty}p_k
    =\sum_{k=0}^{\infty}\frac{\lambda^k}{k!}e^{-\lambda}
    =e^\lambda e^{-\lambda}=1.
\]
\par\endgroup

\begin{theorem}[Теорема Пуассона]
    Рассмотрим серию схем испытаний Бернулли:
    \[
        \xi_n\sim\operatorname{Bin}(n;p_n),\qquad
        np_n\xrightarrow[n\to\infty]{}\lambda,\qquad
        \eta\sim\operatorname{Poiss}(\lambda).
    \]
    Тогда для каждого фиксированного $k=0,1,2,\ldots$
    \[
        \Pr\{\xi_n=k\}\xrightarrow[n\to\infty]{}\Pr\{\eta=k\}.
    \]
\end{theorem}

\begin{Proof}
    \[
        \begin{aligned}
            \Pr\{\xi_n=k\}
            &=C_n^k p_n^k(1-p_n)^{n-k}\\
            &=\frac{n!}{(n-k)!\,k!}\,p_n^k(1-p_n)^{n-k}\\
            &=\frac1{k!}\,
              \underbrace{\frac{n(n-1)\cdots(n-k+1)}{n^k}}
              _{\longrightarrow\,1}\,
              \underbrace{(np_n)^k}_{\longrightarrow\,\lambda^k}\,
              \frac{(1-p_n)^n}
              {\underbrace{(1-p_n)^k}_{\longrightarrow\,1}}.
        \end{aligned}
    \]
    При этом
    \[
        (1-p_n)^n
        =\left(\underbrace{(1-p_n)^{-1/p_n}}
        _{\longrightarrow\,e}\right)^{-np_n}
        \longrightarrow e^{-\lambda}.
    \]
    Следовательно,
    \[
        \Pr\{\xi_n=k\}\longrightarrow
        \frac1{k!}\cdot1\cdot\lambda^k e^{-\lambda}
        =\frac{\lambda^k}{k!}e^{-\lambda}.
    \]
\end{Proof}

\begin{theorem}[Обобщение]
    Пусть $\xi\sim\operatorname{Bin}(n;p)$,
    $\eta\sim\operatorname{Poiss}(\lambda)$, $np=\lambda$.
    Тогда
    \[
        \bigl|\Pr\{\xi\in B\}-\Pr\{\eta\in B\}\bigr|
        \leq np^2
        \qquad\forall B\in\mathcal{B}(\R).
    \]
\end{theorem}

\subsubsection{Геометрическое распределение}

\begin{definition}
    \[
        \operatorname{Geom}(p),\qquad p\in(0;1).
    \]
    Смысл: номер первого успеха в бесконечной схеме испытаний Бернулли.
    \[
        x_i=i\in\N,\qquad p_i=(1-p)^{i-1}p.
    \]
\end{definition}

\begingroup
\small
\setlength{\abovedisplayskip}{6pt}
\setlength{\belowdisplayskip}{6pt}
\noindent Пояснение.
Здесь $p$ --- вероятность успеха в одном испытании,
а $p_i$ --- вероятность первого успеха на $i$-м испытании.
Складываем $p+(1-p)p+(1-p)^2p+\cdots$:
это геометрическая прогрессия со знаменателем $1-p\in(0;1)$;
сумма всех вероятностей равна $1$:
\[
    \sum_{i=1}^{\infty}(1-p)^{i-1}p
    =p\sum_{i=1}^{\infty}(1-p)^{i-1}
    =p\cdot\frac1{1-(1-p)}=1.
\]
\par\endgroup

\subsection{Функция распределения}

\begin{definition}
    \textbf{Функцией распределения случайной величины} $\xi$
    называется функция
    \[
        F_\xi(x)=\Pr\{\xi(\omega)\leq x\},\qquad x\in\R.
    \]
\end{definition}

\begingroup
\tcbset{myboxstyle/.append style={breakable=false}}
\begin{example}
    Пусть $\xi\sim\operatorname{Bern}(p)$, $p=\frac14$.
    \[
        \Pr\{\xi(\omega)\leq0\}
        =\Pr\{\xi(\omega)=0\}=\frac34.
    \]
    \begin{center}
        \begin{tikzpicture}[x=2.5cm,y=3cm,>=Stealth]
            \draw[->] (-0.85,0) -- (2.2,0) node[right] {$x$};
            \draw[->] (0,-0.12) -- (0,1.23) node[above] {$F_\xi(x)$};
            \draw[blue,dashed] (0,1) -- (1,1);
            \draw[blue,dashed] (1,0) -- (1,0.75);
            \draw[blue,very thick] (-0.85,0) -- (0,0);
            \draw[blue,very thick] (0,0.75) -- (1,0.75);
            \draw[blue,very thick] (1,1) -- (2.05,1);
            \filldraw[blue,fill=white] (0,0) circle (1.6pt);
            \filldraw[blue,fill=white] (1,0.75) circle (1.6pt);
            \fill[blue] (0,0.75) circle (1.6pt);
            \fill[blue] (1,1) circle (1.6pt);
            \node[below left] at (0,0) {$0$};
            \node[below] at (1,0) {$1$};
            \node[left] at (0,0.75) {$\frac34$};
            \node[left] at (0,1) {$1$};
        \end{tikzpicture}
    \end{center}
\end{example}
\endgroup

\subsubsection{Свойства функции распределения}

\begin{theorem}
    \begin{enumerate}[label=\arabic*)]
        \item $\lim_{x\to+\infty}F_\xi(x)=1$,
        $\lim_{x\to-\infty}F_\xi(x)=0$.
        \item $F_\xi$ не убывает.
        \item $F_\xi$ непрерывна справа.
    \end{enumerate}
\end{theorem}

\begin{Proof}[Непрерывность справа]
    По Гейне рассмотрим $x_n\downarrow x_0$. Положим
    \[
        A_n=\xi^{-1}\bigl((-\infty;x_n]\bigr),
        \qquad A=\xi^{-1}\bigl((-\infty;x_0]\bigr).
    \]
    Тогда
    \[
        F_\xi(x_n)=\Pr\{\xi(\omega)\leq x_n\}=\Pr(A_n),
        \qquad F_\xi(x_0)=\Pr(A).
    \]
    Поскольку $A_n\downarrow A$, по непрерывности сверху
    \[
        \Pr(A_n)\longrightarrow\Pr(A),
        \qquad F_\xi(x_n)\longrightarrow F_\xi(x_0).
    \]
    Следовательно, $\lim_{x\to x_0+}F_\xi(x)=F_\xi(x_0)$.
\end{Proof}

\begin{remark}
    Для неубывающей, ограниченной и непрерывной справа функции $F$
    на полуинтервалах задаётся
    \[
        m\bigl((a;b]\bigr)=F(b)-F(a).
    \]
\end{remark}

\begin{theorem}[Следствие из теоремы Каратеодори]
    Существует взаимно однозначное соответствие:
    \[
        \text{вероятностные меры на }\mathcal{B}(\R)
        \quad\longleftrightarrow\quad
        \text{функции распределения}.
    \]
\end{theorem}

\subsection{Абсолютно непрерывное распределение}

\begin{definition}
    Распределение называется \textbf{абсолютно непрерывным},
    если у него есть плотность $p_\xi$, то есть
    \[
        p_\xi(x)\geq0,\qquad
        F_\xi(x)=\int_{-\infty}^{x}p_\xi(t)\,\diff t.
    \]
\end{definition}

\subsection{Сравнение случайных величин}

\begin{exercise}
    Пусть $\xi,\eta$ --- случайные величины. Тогда
    \[
        \{\omega\colon\xi(\omega)<\eta(\omega)\}\in\mathcal{F}.
    \]
\end{exercise}

\begin{solution}
    Пусть $\Q=\{x_n\}_{n\in\N}$. Тогда
    \[
        \begin{aligned}
            &\{\omega\colon\xi(\omega)<\eta(\omega)\}\\
            &\quad=\bigcup_{n=1}^{\infty}
            \left(\underbrace{\{\omega\colon\xi(\omega)<x_n\}}
            _{\in\mathcal{F}}\cap
            \underbrace{\{\omega\colon x_n<\eta(\omega)\}}
            _{\in\mathcal{F}}\right)\in\mathcal{F}.
        \end{aligned}
    \]
\end{solution}
